% Template para uso no TCC do Departamento de F�sica de Ji-Paran�
% da Universidade Federal de Rond�nia
% Elaborado em 10 de junho de 2017 por Marco Polo Moreno de Souza com base no template book.
% Aperfei�oado em 07 de mar�o de 2019, com a introdu��o de macros
% Modifica��o na estrutura do de oneside para twoside 1 de setembro de 2019
%Resolu��o de warning.

\documentclass[a4paper,reqno,12pt,twoside]{book}
\small
\RequirePackage{fix-cm}
\usepackage{amsmath}%
\usepackage{amsfonts}%
\usepackage{amssymb}%
%\usepackage{graphicx}%
\usepackage{physics}%
\usepackage{blindtext}
\usepackage
[ruled, % Caixa em volta, leg. acima, linha ap�s leg.
portuguese] % Para algoritmos em portugu�s.
{algorithm2e} % Enfim o pacote.
%Veja o arquivo abaixo antes de adicionar outras bibliotecas
\usepackage[mode=buildnew]{standalone}% requires -shell-escape


\usepackage{calc}
%\usepackage{svg}
\usepackage{pst-text}
\newcommand{\executeiffilenewer}[3]{%
	\ifnum\pdfstrcmp{\pdffilemoddate{#1}}%
	{\pdffilemoddate{#2}}>0%
	{\immediate\write18{#3}}\fi%
}
\newcommand{\includesvg}[1]{%
	\executeiffilenewer{#1.svg}{#1.pdf}%
	{inkscape -z -D --file=#1.svg %
		--export-pdf=#1.pdf --export-latex}%
	\input{#1.pdf_tex}%
}
\usepackage{cancel}
\usepackage{scrextend}
\usepackage{indentfirst}
\usepackage{float}
\usepackage[brazilian]{babel}
\usepackage[latin1]{inputenc}
\usepackage{times}
\usepackage[T1]{fontenc}
\usepackage[final]{pdfpages}
\usepackage{epstopdf}
\usepackage[top=3cm,bottom=2cm,left=3cm,right=2cm,marginparsep=0pt]{geometry}
\usepackage{multirow}
\usepackage[colorlinks=true,linkcolor=blue,citecolor=red,urlcolor=black,bookmarks=true,pdfstartview=FitB]{hyperref} %uso de links
%\usepackage[hyphenbreaks]{breakurl}
%\usepackage{makeidx}
%\makeindex
\usepackage{wrapfig}
\usepackage{flashmovie}
\usepackage{fancyhdr} %cabe�alho
\usepackage{setspace} %espa�amento entre linhas

\usepackage{titlesec}    
\titleformat{\chapter}[display]
{\normalfont%
	\Large% %change this size to your needs for the first line
	\bfseries}{\chaptertitlename\ \thechapter}{20pt}{%
	\Large %change this size to your needs for the second line
}

\usepackage{xcolor}
\definecolor{verde}{rgb}{0,0.5,0}
\usepackage{showexpl} % exibe codigos de figuras
\usepackage{listings}
\lstset{
	language=C++,
	basicstyle=\ttfamily\small, 
	keywordstyle=\color{blue}, 
	stringstyle=\color{verde}, 
	commentstyle=\color{red}, 
	extendedchars=true, 
	showspaces=false, 
	showstringspaces=false, 
	numbers=left,
	numberstyle=\tiny,
	breaklines=true, 
	backgroundcolor=\color{green!10}, 
	breakautoindent=true, 
	captionpos=b,
	xleftmargin=0pt,
	emph={main, printf, scanf},
	emphstyle={\color{black}\bf},
	morekeywords={comando1, comando2},language=C++,
	basicstyle=\ttfamily\small,
	keywordstyle=\color{blue},
	stringstyle=\color{verde},
	commentstyle=\color{red},
	extendedchars=true,
	showspaces=false,
	showstringspaces=false,
	numbers=left,
	numberstyle=\tiny,
	breaklines=true,
	backgroundcolor=\color{green!10},
	breakautoindent=true,
	captionpos=b,
	xleftmargin=0pt,
}
\usepackage{amsmath}
\usepackage[font=footnotesize, labelfont=bf, margin=0.5cm]{caption} %Altera formata��o das legendas
\usepackage{indentfirst} %Espa�amento aplicado � primeira linha do primeiro par�grafo
%%%%%%%%%%%%%%%%%%%%%%
\usepackage{color}
\usepackage{ifpdf}
\ifpdf %if using pdfLaTeX in PDF mode
%\usepackage[pdftex]{graphicx}
\DeclareGraphicsExtensions{.pdf,.png,.jpg,.jpeg,.mps}
\usepackage{pgf}
\usepackage{tikz}
\else %if using LaTeX or pdfLaTeX in DVI mode
\usepackage{graphicx}
\DeclareGraphicsExtensions{.eps,.bmp}
\DeclareGraphicsRule{.emf}{bmp}{}{}% declare EMF filename extension
\DeclareGraphicsRule{.png}{bmp}{}{}% declare PNG filename extension
\usepackage{pgf}
\usepackage{tikz}
\usepackage{pstricks}
\fi
\usepackage{epic,bez123}
\usepackage{floatflt}% package for floatingfigure environment
\usepackage{wrapfig}% package for wrapfigure environment

%%%%%%%%%%%%%%%

\makeatletter
\newcommand{\showfont}
{%
	(encoding: \f@encoding{},
	family: \f@family{},
	series: \f@series{},
	shape: \f@shape{},
	size: \f@size{},
	baseline: \f@baselineskip{})
	%tfsize: \tf@size{},
	%sfsize: \sf@size{},
	%sssize: \ssf@size{}
}
\makeatother

%%%%%%%%%%%%%%%%%%%%%%%%%%%%%%%%%%%%%%%%%%%%%%%%%%%%%%%%%%%%%%%

%Preencha os campos abaixo conforme seu TCC
\newcommand{\Titulo}{Propaga��o de pulsos de luz em sistemas at�micos}
\newcommand{\Nome}{Eliton Trindade Gomes}
\newcommand{\curso}{Bacharelado} %Digite Licenciatura Plena ou Bacharelado
\newcommand{\graduacao}{Bacharel} %Digite Licenciado ou Bacharel conforme o caso
\newcommand{\orientador}{Dr. Marco Polo Moreno de Souza} %Nome com t�tulo. Ex: Dr. Marco Polo Moreno de Souza
\newcommand{\primeiromembro}{Nome do professor da banca} %Nome com t�tulo. Ex: Dr. Fulano de Tal
\newcommand{\segundomembro}{Nome do professor da banca} %Nome com t�tulo. Ex: MSc. Beltrano
\newcommand{\instituicao}{DEFIJI/CJP/UNIR} %Institui��o do segundo membro. Se for interno, digite DEFIJI/CJP/UNIR
\newcommand{\Cidade}{Ji-Paran�, RO}
\newcommand{\Data}{M�s e ano da defesa}

%%%%%%%%%%%%%%%%%%%%%%%%%%%%%%%%%%%%%%%%%%%%%%%%%%%%%%%%%%%%%%%
%Dados da defesa, a ser preenchido pelo orientador
\newcommand{\dia}{xxx} %Ex: quinze
\newcommand{\mes}{xxx} %Ex: abril
\newcommand{\ano}{xxx} %Ex: 2019
\newcommand{\inicio}{xxx} %In�cio da defesa. Ex: 10:15 h
\newcommand{\fim}{xxx} %Fim da sess�o. Ex: 11:45 h
\newcommand{\arguicao}{xxx} %Dura��o da argui��o da banca em minutos. Ex: 45
\newcommand{\arguicaoextenso}{xxx} %Idem acima, por extenso. Ex: quarenta e cinco
\newcommand{\local}{xxx} %Ex: audit�rio do Campus da UNIR de Ji-Paran�
\newcommand{\nota}{xxx} %Ex: 8,5
\newcommand{\notaextenso}{xxx} %Ex: oito inteiros e cinco d�cimos

\newcommand{\titulo}{\expandafter\MakeUppercase\expandafter{\Titulo}}
\newcommand{\nome}{\expandafter\MakeUppercase\expandafter{\Nome}}
\newcommand{\cidade}{\expandafter\MakeUppercase\expandafter{\Cidade}}
\newcommand{\data}{\expandafter\MakeUppercase\expandafter{\Data}}
\titleformat{\chapter}
{\normalfont\bfseries}{\thechapter}{0.3em}{\MakeUppercase}
\titlespacing*{\chapter}{0pt}{-36pt}{36pt}

\titleformat{\section}
{\normalfont}{\thesection}{0.3em}{\MakeUppercase}
\titlespacing*{\section}{0pt}{36pt}{36pt}

\titleformat{\subsection}
{\normalfont\bfseries}{\thesubsection}{0.3em}{}
\titlespacing*{\subsection}{0pt}{36pt}{36pt}

\setlength{\parindent}{1.25cm} %Primeira linha do par�grafo
%\renewcommand{\baselinestretch}{1.5} %
%\setstretch{1.5}

%Coloca a numera��o no canto externo da p�gina (necess�rio para impress�o frente e verso)
\pagestyle{fancy}
\fancyhf{}
\renewcommand{\headrulewidth}{0pt}% elimina linhas horizontais
\fancyhead[RO, LE]{\thepage}

%Faz o mesmo que a sequ�ncia acima em p�ginas onde se come�am cap�tulos
\fancypagestyle{plain}{%
  \renewcommand{\headrulewidth}{0pt}%
  \fancyhf{}%
  \fancyhead[RO, LE]{\thepage}
}

\begin{document}
	\pagenumbering{Alph}
	\renewcommand{\bibname}{Refer�ncias}
	\chapter*{}
\thispagestyle{empty}

\vspace{-1cm}
\begin{center}
\textbf{\nome}
\end{center}

\vspace{3.5cm}
\begin{figure}[htbp]
	\centering
		\includegraphics[width=0.99\textwidth]{figuras_do_template/logo.eps}
\end{figure}
\vspace{4cm}

\begin{center}
\textbf{\titulo}
\end{center}

\vspace{8cm}
\begin{center}
\textbf{\cidade} \\ \textbf{\data}
\end{center}
	\chapter*{}
\thispagestyle{empty}

\vspace{-1cm}

\begin{center}
\textbf{\nome}
\end{center}

\vspace{11cm}

\begin{center}
\textbf{\titulo}
\end{center}

\vspace{1cm}

\begin{minipage}{0.4\linewidth}
.
\end{minipage}
\begin{minipage}{0.5\linewidth}
Trabalho de Conclus�o de Curso apresentado
ao Departamento de F�sica de Ji-Paran�,
Universidade Federal de Rond�nia, Campus de
Ji-Paran�, como parte dos quesitos para a
obten��o do T�tulo de {\graduacao} em F�sica, sob orienta��o do Prof.
{\orientador}.
\end{minipage}

\vspace{5cm}

\begin{center}
\textbf{\cidade}\\
\textbf{\data}
\end{center}
	\chapter*{}
\thispagestyle{empty}
\vspace{-2cm}

\begin{figure}[htbp]
	\centering
		\includegraphics[width=0.95\textwidth]{logo-ata.eps}
\end{figure}
\vspace{0.3cm}
\begin{center}
\textbf{ATA DE AVALIA��O DO TRABALHO DE CONCLUS�O DE CURSO DO CURSO DE (LICENCIATURA PLENA/BACHARELADO) EM F�SICA}
\end{center}
\vspace{0.7cm}
\doublespacing

\noindent Aos {\dia} dias do m�s de {\mes} do ano de {\ano}, �s {\inicio}, no {\local}, reuniu-se a Banca Julgadora composta pelo professor orientador {\orientador} e pelos examinadores {\primeiromembro} e {\segundomembro}, para avaliarem o Trabalho de Conclus�o de Curso, do Curso de
{\curso} em F�sica, intitulado ``\textbf{\titulo}'', do discente {\textit{\nome}}. Ap�s a apresenta��o, o candidato foi arguido pelos integrantes da Banca Julgadora por {\arguicao} ({\arguicaoextenso}) minutos. Ao final da argui��o, a Banca Julgadora, em sess�o reservada, \textbf{aprovou} o candidato com nota {\nota} ({\notaextenso}), em uma avalia��o de 0 (zero) a 10 (dez). Nada mais havendo a tratar, a sess�o foi encerrada �s {\fim}, dela sendo lavrada a presente ata, assinada por todos os membros da Banca Julgadora.

\vspace{2cm}
\singlespacing

\begin{center}
\rule{13cm}{0.1mm}\\
Prof. {\orientador} - DEFIJI/CJP/UNIR \\
Orientador
\end{center}

\vspace{1cm}
\begin{center}
\rule{13cm}{0.1mm}\\
Prof. {\primeiromembro} - DEFIJI/CJP/UNIR
\end{center}

\vspace{1cm}
\begin{center}
\rule{13cm}{0.1mm}\\
Prof. {\segundomembro} - \instituicao
\end{center}
	%%%%%%%%%%%%%%%%%%%%%%%%%%%%%%%%%%%%%%%%%%%%%%%%%%%%%%%%%%%%%%%
\chapter*{Dedicat�ria}
\thispagestyle{empty}
%A organiza��o das p�ginas de dedicat�ria, agradecimentos e ep�grafe ficam a crit�rio do aluno, sendo opcionais.

Digite a dedicat�ria aqui.


%%%%%%%%%%%%%%%%%%%%%%%%%%%%%%%%%%%%%%%%%%%%%%%%%%%%%%%%%%%%%%%
	%%%%%%%%%%%%%%%%%%%%%%%%%%%%%%%%%%%%%%%%%%%%%%%%%%%%%%%%%%%%%%%

\chapter*{Agradecimentos}
\thispagestyle{empty}

Digite os agradecimentos aqui.


%%%%%%%%%%%%%%%%%%%%%%%%%%%%%%%%%%%%%%%%%%%%%%%%%%%%%%%%%%%%%%%
	%%%%%%%%%%%%%%%%%%%%%%%%%%%%%%%%%%%%%%%%%%%%%%%%%%%%%%%%%%%%%%%
\chapter*{Ep�grafe}
\thispagestyle{empty}

Digite a ep�grafe aqui.


%%%%%%%%%%%%%%%%%%%%%%%%%%%%%%%%%%%%%%%%%%%%%%%%%%%%%%%%%%%%%%%
	%%%%%%%%%%%%%%%%%%%%%%%%%%%%%%
\chapter*{Resumo}
\thispagestyle{empty}
\noindent
O resumo em l�ngua portuguesa dever� conter no m�nimo 150 e no m�ximo 500 palavras.
Bla Bla Bla Bla Bla Bla Bla Bla Bla Bla Bla Bla Bla Bla Bla Bla Bla Bla 
Bla Bla Bla Bla Bla Bla Bla Bla Bla Bla Bla Bla Bla Bla Bla Bla Bla Bla 
Bla Bla Bla Bla Bla Bla Bla Bla Bla Bla Bla Bla Bla Bla Bla Bla Bla Bla 
Bla Bla Bla Bla Bla Bla Bla Bla Bla Bla Bla Bla Bla Bla Bla Bla Bla Bla 
Bla Bla Bla Bla Bla Bla Bla Bla Bla Bla Bla Bla Bla Bla Bla Bla Bla Bla 
Bla Bla Bla Bla Bla Bla Bla Bla Bla Bla Bla Bla Bla Bla Bla Bla Bla Bla 
Bla Bla Bla Bla Bla Bla Bla Bla Bla Bla Bla Bla Bla Bla Bla Bla Bla Bla 
Bla Bla Bla Bla Bla Bla Bla Bla Bla Bla Bla Bla Bla Bla Bla Bla Bla Bla 
Bla Bla Bla Bla Bla Bla Bla Bla Bla Bla Bla Bla Bla Bla Bla Bla Bla Bla 
Bla Bla Bla Bla Bla Bla Bla Bla Bla Bla Bla Bla Bla Bla Bla Bla Bla Bla 
Bla Bla Bla Bla Bla Bla Bla Bla Bla Bla Bla Bla Bla Bla Bla Bla Bla Bla 
Bla Bla Bla Bla Bla Bla Bla Bla Bla Bla Bla Bla Bla Bla Bla Bla Bla Bla 
Bla Bla Bla Bla Bla Bla Bla Bla Bla Bla Bla Bla Bla Bla Bla Bla Bla Bla 
Bla Bla Bla Bla Bla Bla Bla Bla Bla Bla Bla Bla Bla Bla Bla Bla Bla Bla 
Bla Bla Bla Bla Bla Bla Bla Bla Bla Bla Bla Bla Bla Bla Bla Bla Bla Bla 
Bla Bla Bla Bla Bla Bla Bla Bla Bla Bla Bla Bla Bla Bla Bla Bla Bla Bla 
Bla Bla Bla Bla Bla Bla Bla Bla Bla Bla Bla Bla Bla Bla Bla Bla Bla Bla 
Bla Bla Bla Bla Bla Bla Bla Bla Bla Bla Bla Bla Bla Bla Bla Bla Bla Bla 

\vspace{1.0cm}
\noindent
\textbf{Palavras-chave}: palavra-chave 1. palavra-chave 2. palavra-chave 3.
% Minimo de 3 e m�ximo de 6 palavras-chave.

\listoftables
\thispagestyle{empty}
\listoffigures
\thispagestyle{empty}

\tableofcontents
\thispagestyle{empty}

\mainmatter


%%%%%%%%%%%%%%%%%%%%%%%%%%%%%%%%%%%%%%%%%%%%%%%%%%%%%%%%%%%%%%%
	%%%%%%%%%%%%%%%%%%%%%%%%%%%%%%%%%%%%%%%%%%%%%%%%%%%%%%%%%%%%%%%
\chapter{Introdu��o}
\label{introducao}

A inten��o deste TCC � estudar a propaga��o de um pulso de luz ultracurto por uma  amostra de muito �tomos, onde cada �tomo pode ser aproximado como um sistema de dois n�veis. Para isso, modelamos nosso problema usando �tomos de Rub�dio-87 e um pulso ressonante com a frequ�ncia da transi��o  $5^2S_{1/2}\longrightarrow5^2P_{1/2}$ desse �tomo.

Dessa forma, resolvemos numericamente as equa��es de Maxwell-bloch. Nesta equa��es s�o acopladas as equa��es de Bloch, que descrevem a evolu��o temporal da matriz densidade conforme o pulso provoca transi��es do estado fundamental  $5^2S_{1/2}$ para o estado excitado $5^2P_{1/2}$, produzindo, assim, uma polariza��o no meio. Tamb�m est� acoplada uma aproxima��o das equa��es de Maxwell que descrevem as modifica��es no pulso devido essa polariza��o.
 
Este trabalho esta organizado da seguinte forma:

No Cap�tulo \ref{cap:capitulo3}, apresentamos um resumo da leis de Maxwell e obtemos uma aproxima��o para propaga��o de pulso em meios lineares e n�o lineares, separando o pulso e uma parte lenta (termo de modula��o) que oscila rapidamente (termo de propaga��o).    

No Cap�tulo \ref{cap:capitulo2}, trazemos o formalismo de matriz densidade que utilizamos pra tratar o ensembles de sistemas de dois n�veis. Mostrando como obter esse formalismo a partir do formalismo usual da mec�nica qu�ntica (auto valores e auto estados). Mostramos tamb�m as principais propriedades da matriz densidade e como descrever sua evolu��o temporal.

No cap�tulo \ref{cap:capitulo4}, discutimos a intera��o de campo eletromagn�tico com sistemas de dois n�veis em um ponto de vista semicl�ssico, onde o campo � tratado classicamente e o sistema de dois n�veis e tratado pela mec�nica qu�ntica. Pra isso, mostramos como � obtido o hamiltoniano de aproxima��o dipolar, como obtemos a equa��es �ptica de bloch e aplicamos para uma intera��o de um campo cont�nuo com �tomos de Rub�dio-87, onde estudamos a transi��o $5^2S_{1/2}\longrightarrow5^2P_{1/2}$, considerando casos em que o campo � ressonante ou pr�ximo da resson�ncia at�mica, regime onde as equa��es �pticas de Bloch s�o validas. 

Por fim, o Cap�tulo \ref{cap:maxwellbloch} traz o estudo da propaga��o de pulso ultracurto em sistema de muitos �tomos. Para isso, obtemos as equa��es de  maxwell-bloch e apresentamos o teorema da �rea. Ao final do Cap�tulo trazemos e discutimos sobre  os resultados num�ricos para essa equa��es de Maxwell-Bloch, para um pulso ultracurto interagindo com uma amostra de �tomos de Rub�dio-87, onde pulso � ressonante com a transi��o~$5^2S_{1/2}\longrightarrow5^2P_{1/2}$.      
  


%%%%%%%%%%%%%%%%%%%%%%%%%%%%%%%%%%%%%%%%%%%%%%%%%%%%%%%%%%%%%%%
	
%%%%%%%%%%%%%%%%%%%%%%%%%%%%%%%%%%%%%%%%%%%%%%%%%%%%%%%%%%%%%%%
\chapter{Mec�nica Qu�ntica e Operador Densidade ($\rho$) }
\label{capitulo2}

%Para cada cap�tulo adicional, digite o comando \chapter{T�tulo do Cap�tulo}

Neste capitulo nos dedicamos a discutir sobre formalismo do operador densidade , desenvolvido por J.von Neumann em 1917, e suas vantagens em rela��o ao formalismos de autoestados e autovetores no estudo de sistema qu�nticos.

\section{Matriz densidade}
\label{secao2.1}

Como sabemos, o formalismo usual da mec�nica qu�ntica (formalismo de dirac) nos permite fazer previs�es sobre um conjunto de sistemas f�sicos elaborados de forma id�ntica (ensemble puro). Em termos mais especifico, garantimos que todos o sistemas membros deste ensemble seja caracterizado por um mesmo ket de estados $\ket{\alpha}$. Assim, este formalismo n�o � v�lido se considerarmos, por exemplo, que $70\%$ desses sistemas s�o caracterizados pelo ket de estado $\ket{\alpha}$ e $30\%$ pelo ket de estado $\ket{\beta}$(ensemble misto). Para lidar com essa situa��o, precisamos introduzir o conceito de operador densidade, que nos permitir� descrever quantitativamente conjuntos de sistemas qu�ntico para ensemble puros ou, at� mesmo, ensemble mistos completamente aleat�rios. 

Consideremos o ensemble misto, onde uma fra��o de membros com popula��o relativa $\alpha^{(1)}$ � caracterizada por $\ket{\alpha^{(1)}}$; outra fra��o $\alpha^{(^2)}$ � caracterizada por $\ket{\alpha^{(2)}}$ e assim sucessivamente. Podemos dizer, com certa precis�o, que um ensemble misto pode ser visto como uma mistura de ensembles puros. As popula��es $\alpha^{(i)}$ devem satisfazer a condi��o de normaliza��o, ou seja:

\begin{equation} \label{eq:1}
	\sum_i \alpha^{(i)} = 1 
\end{equation}
   
N�o � necess�rio que $\alpha^{(1)}$, $\alpha^{(2)}$,..., $\alpha^{(i)}$ sejam ortogonais entre si e o n�mero de termos na soma em $i$ \eqref{eq:1} n�o precisam igual ao numero de dimens�es $N$ no espa�o de Hilbert.
Vamos supor que realizamos a medida de um observ�vel A em um ensemble misto. � poss�vel calcular o valor m�dio se houver um numero grande de medidas. O resultado � dado pela m�dia sobre o ensembles, definida por:



\subsection{Subse��o}

Digite o conte�do do uma subse��o aqui.

\section{Figuras}

Exemplo de figura:

\begin{figure}[H]
	\centering
	\includegraphics[width=0.6\linewidth]{figuras_do_trabalho/fig1.eps}
	\caption{ggggggggggggggggg}
	\label{fig1}
\end{figure}



%Cita��o de figura: \ref{fig1}.

\section{Equa��es}

Exemplo de equa��o centralizada:

\begin{equation}
\label{pitagoras}
	a^2 = b^2 + c^2.
\end{equation}

Exemplo de equa��o no texto: $e^{ix} = \cos{x} + i\sin{x}$. Cita��o de equa��o: \ref{pitagoras}.

\section{Tabelas}

Exemplo de tabela:

\begin{table}[H]
\begin{center}
	\begin{tabular}{|c|c|c|}
		\hline
		$\sin{x}$ & $\cos{x}$ & $\tan{x}$ \\
		\hline
		$\sec{x}$ & $\csc{x}$ & $\cot{x}$ \\
		\hline
		$\arcsin{x}$ & $\arccos{x}$ & $\arctan{x}$ \\
		\hline
		$\sinh{x}$ & $\cosh{x}$ & $\tanh{x}$ \\
		\hline
	\end{tabular}
	\vspace{0.3cm}
	\caption{Fun��es trigonom�tricas e hiperb�licas.}
	\label{funcoes}
\end{center}
\end{table}

%Cita��o de tabela: \ref{funcoes}

\section{C�digos}

Exemplo de c�digo (linguagem C):

\vspace{0.5cm}
\begin{lstlisting}
#include<stdio.h>

int k;

main()
{
	for (k=1; k<=5; k++)
		printf("F�sica - UNIR - Ji-Paran�\n");
}
\end{lstlisting}

\section{Cita��o}

Exemplo de cita��o:

\vspace{0.5cm}
\begin{addmargin}[3cm]{0.0cm}\small{
		Bla Bla Bla Bla Bla Bla Bla Bla Bla Bla Bla Bla Bla Bla Bla Bla Bla Bla Bla Bla
		Bla Bla Bla Bla Bla Bla Bla Bla Bla Bla Bla Bla Bla Bla Bla Bla Bla Bla Bla Bla
		Bla Bla Bla Bla Bla Bla Bla Bla Bla Bla Bla Bla Bla Bla Bla Bla Bla Bla Bla Bla
		Bla Bla Bla Bla Bla Bla Bla Bla Bla Bla Bla Bla Bla Bla Bla Bla Bla Bla Bla Bla
		Bla Bla Bla Bla Bla Bla Bla Bla Bla Bla Bla Bla Bla Bla Bla Bla Bla Bla Bla Bla
		Bla Bla Bla Bla Bla Bla Bla Bla Bla Bla Bla Bla Bla Bla Bla Bla Bla Bla Bla Bla
		Bla Bla Bla Bla Bla Bla Bla Bla Bla Bla Bla Bla Bla Bla Bla Bla Bla Bla Bla Bla
}\end{addmargin}
\vspace{0.5cm}

Citando um trabalho: \cite{ARAUJO} ou (ARA�JO, 2004).


%%%%%%%%%%%%%%%%%%%%%%%%%%%%%%%%%%%%%%%%%%%%%%%%%%%%%%%%%%%%%%%
\chapter[Outro cap�tulo]{Outro cap�tulo}

Digite aqui o conte�do de outro cap�tulo.

\begin{table}[H]
\begin{center}
	\begin{tabular}{|c|c|c|}
		\hline
		$\sin{x}$ & $\cos{x}$ & $\tan{x}$ \\
		\hline
		$\sec{x}$ & $\csc{x}$ & $\cot{x}$ \\
		\hline
		$\arcsin{x}$ & $\arccos{x}$ & $\arctan{x}$ \\
		\hline
		$\sinh{x}$ & $\cosh{x}$ & $\tanh{x}$ \\
		\hline
	\end{tabular}
	\vspace{0.3cm}
	\caption{Fun��es trigonom�tricas e hiperb�licas.}
	\label{funcoes2}
\end{center}
\end{table}

\begin{figure}[H]
	\centering
	\includegraphics[width=0.6\linewidth]{figuras_do_trabalho/fig1.eps}
	\label{fig2}
	\caption{Espectro de um laser de femtossegundos.}
\end{figure}


%%%%%%%%%%%%%%%%%%%%%%%%%%%%%%%%%%%%%%%%%%%%%%%%%%%%%%%%%%%%%%%
	%%%%%%%%%%%%%%%%%%%%%%%%%%%%%%%%%%%%%%%%%%%%%%%%%%%%%%%%%%%%%%%

\chapter{Conclus�o}

Digite a conclus�o do TCC aqui.


%%%%%%%%%%%%%%%%%%%%%%%%%%%%%%%%%%%%%%%%%%%%%%%%%%%%%%%%%%%%%%%
	%Digite as refer�ncias abaixo
	\bibliography{referencias.bib}
	\bibliographystyle{ieeetr}
	
	
	\backmatter \appendix
%%%%%%%%%%%%%%%%%%%%%%%%%%%%%%%%%%%%%%%%%%%%%%%%%%%%%%%%%%%%%%%
\chapter{T�tulo do Primeiro Ap�ndice}

Digite o primeiro ap�ndice aqui.


%%%%%%%%%%%%%%%%%%%%%%%%%%%%%%%%%%%%%%%%%%%%%%%%%%%%%%%%%%%%%%%
\chapter{T�tulo do Segundo Ap�ndice}
%Para cada ap�ndice adicional, insira o comando \chapter{T�tulo do Ap�ndice}

Digite o segundo ap�ndice aqui.
	
\end{document}